\documentstyle{report}
\title{User's Manual and Programming Guide for THOR/HORSE \\ Version 1.0}
\author{Klaus J. Berkling \\ Michael L. Hilton \\ \\ School of Computer and Information Sciences \\
Syracuse University\\Syracuse, NY 13244  USA \\ \\
\copyright 1990}

\begin{document}
\maketitle
\tableofcontents

\chapter{Introduction}

	This document serves two functions ---  first, it is a
programmer's guide to The Head Order Reduction Language
(THOR), a lazy, untyped functional language with normal-order semantics;
second, it is a
user's manual for the Head Order Reduction System (HORSE), which
is a programming environment for  THOR.  This document is not a
tutorial; although the material presented is not technically
difficult, it assumes that the reader has some familiarity
with the principles and lexicon of the lambda calculus and functional
programming.


\section{THOR}

	THOR is an implementation of an enriched lambda calculus. 
The pure lambda calculus has a
simple syntax and semantics which deals only with expressions
composed of variables, applications, and abstractions.  While
these are all that is needed to compute all computable functions,
it is more convenient to `enrich' the pure lambda calculus with
other kinds of objects, like integers and floating point numbers,
and primitive operations to manipulate these objects.

	THOR is intended to be used as the assembly language of the
RED 1.5 reduction machine, which is currently being designed at
Syracuse University.  THOR is in itself a functional
programming language, albeit a rather limited old-fashioned one when
compared to the `modern' functional languages (such as
Haskell) for which THOR is designed to be the compilation target language.
  When faced with it's deficiencies (such as lack of input-output
primitives), it is helpful to remember that THOR is still under
development, and is not yet an industrial strength language.
The syntax and semantics of THOR are presented in 
Chapter~\ref{chapter-THOR}.

	THOR is similar in many ways to early functional languages such as 
Henderson's LISPKIT-LISP \cite{henderson}.  There are, however, 
a few important exceptions:
\begin{itemize}
	\item THOR is a {\em reduction\/} language.
	\item THOR reduces expression to {\em strong head normal form\/}.
	\item Computation in THOR is controllable by the user.
\end{itemize}
It is these exceptions that combine to make THOR unique among current
programming languages.  Each of these points will now be expanded
further.


\subsection{Reduction Languages}

	In a reduction language, computation is accomplished through
equivalence preserving transformations performed on expressions. 
In general, these transformations make an expression simpler,
hence the name `reduction' language.  
	\footnote{It should be noted that being a reduction language
	generally has more to do with how a language is implemented
	than with any property of the language itself.  For
	instance, it is possible to implement pure LISP or Scheme as
	a reduction language.}
The process of transforming an expression into an equivalent,
simpler one is usually referred to as {\em reduction\/}.

	Because all computation proceeds via equivalence preserving
transformations, reduction languages have
semantics that are relatively easy to reason about.  That is
a big plus when it comes to proving properties of programs.  

	All programming languages can be viewed as transformation
systems in the sense that a body of code is transformed, via some
mechanism, into a result of some sort.  With most conventional
programming languages, this transformation system accepts code as
input and yields some sort of data objects (which are created by
executing this code) as output.  Thus,
there is a fundamental difference in the input and output types
of the transformation system.  This limits the ability of
software to transform itself without using some sort of external
meta-programming mechanism.

	Reduction languages, on the other hand, transform
expressions into expressions, and thus the input and output types
are the same.  It is therefore possible for expressions to
transform themselves in useful ways, such as producing a
specialized version of itself which is parameterized to specific
input data.  The idea of software transforming itself seems very
radical and is a powerful feature of THOR.



\subsection{Strong Head Normalization}

	Perhaps the most striking difference between THOR and other
current functional languages is that THOR reduces expressions fully to
strong head normal form,
	 \footnote{This should not be confused with {\em strong normalization},
	  which is often used to describe systems in which all expressions have
 	  a normal form.}
whereas most other languages only evaluate
expressions to weak head normal form.  Strong head normalizing 
means that the body of an abstraction will be
reduced even if the abstraction is not applied to an argument.
For example, in a strong head normalizing language, the expression
\( (\lambda x . (+ \ x \  (* \ 2 \  3))) \) will be reduced to 
\( (\lambda x . (+ \ x \ 6)) \), whereas in a weak head normalizing
language, the expression will be left in its original
form.

	Strong head normalization is uncommon in programming languages for several
reasons:
\begin{itemize}
	\item Many language designers fail to realize it is a useful feature.
	\item For most functional language implementation, stopping reduction
			when weak head	normal form is reached is essential to the lazy 
			evaluation of data structures. 
			THOR handles lazy data structures in a way that does not
			depend on stopping reduction when weak head normal form is reached.
	\item Strong head normalization is not really useful unless there is a
			mechanism for reconstructing the form of reduced expressions.  
			Most languages do not provide such a mechanism; THOR does.  
	\item Reducing to strong head normal form implies a language must
			have some way to deal with name clash problems (see
			Section~\ref{sec-protect}).  
\end{itemize}


\subsection{Controlled Reduction}

	Each equivalence preserving transformation in THOR is
semantically a single atomic action.  Through the THOR
programming environment, HORSE, the user is given the opportunity
to specify how many transformations to perform.  This allows a
user to step through a computation and observe how an expression
behaves, or to generate an intermediate expression which can be
put to various uses.


\section{HORSE}

	HORSE is the programming environment for THOR.  It allows
the user to do things like assign names to expressions, load in
files containing expressions, and control the reduction process
in various ways.  HORSE is described in detail in Chapter~\ref{chapter-horse}.


\include{horse}
\include{thor}
\include{example}

\begin{thebibliography}{1}

\bibitem{berkling} Klaus Berkling. ``Head Order Reduction:  A Graph
Reduction Scheme for the Operational Lambda Calculus,'' {\em Graph
Reduction}, Springer-Verlag Lecture Notes In Computer Science, Vol. 279,
pp. 26--48, 1987.

\bibitem{henderson} Peter Henderson. {\em Functional Programming: 
Application and Implementation.}  Prentice-Hall, 1980. 

\bibitem{roylance} Gerald Roylance.  ``Expressing Mathematical
Subroutines Constructively,'' {\em Proceedings of the 1988 ACM Conference on
Lisp and Functional Programming}, July 25-27 1988, Snowbird, Utah, pp. 8--13.

\end{thebibliography}

\end{document}
